\section{Introduction}
Recovering a positive distribution from a small number of Laplace-transform samples is a prototypical ill-conditioned inverse problem. Diffusion-ordered nuclear magnetic resonance spectroscopy (DOSY) provides a particularly useful setting: the acquisition dimension contains attenuations, while the spectral dimension contains many potentially related signals \cite{stejskal,morris}. Different chemical shifts can carry contributions with the same diffusion coefficient, so independent inversion of every frequency discards useful structure. Conversely, imposing excessive sharing can spread a fitted component into spectral regions where it is absent.

Two questions must therefore be separated. The first is how well a positive measure explains the observations. The second is which components, if any, are sufficiently resolved to support an interpretation. A small residual answers neither the second question nor the question of whether a distribution is atomic or continuous. A sharply drawn DOSY band can be the consequence of the chosen representation, and a smooth band can be the consequence of regularization. Neither appearance alone establishes molecular identity.

Classical approaches address different aspects of Laplace inversion. CONTIN uses constrained regularization \cite{contin}; TRAIn uses an iterative trust-region construction \cite{train}; ITAMeD applies iterative thresholding \cite{itamed}; and PALMA combines entropy and sparsity \cite{palma}. Regularized Kaczmarz methods provide an algebraic reconstruction perspective \cite{kaczmarz}. Joint processing is also established: simultaneous inversion with nonnegative low-rank regularization was investigated by Yuan et al.\ \cite{silt}, and high-resolution DOSY reconstruction by Lin et al.\ \cite{lin}. These precedents rule out treating joint inversion, positivity, or a low-rank factorization as novel in themselves.

The diffusion-NMR work of Arrabal-Campos and collaborators supplies specific application and algorithmic precedents. Solvent-viscosity-aware molecular-weight prediction, together with its published correction, establishes the need to separate diffusion estimation from molecular-weight calibration \cite{arrabal2016,arrabal2016correction}. Subsequent protein and polymer studies address different viscosity and concentration settings \cite{protein,concentration}; these applications motivate, but do not validate, the current synthetic reconstruction study. Time-resolved diffusion NMR has also been used to monitor photoinduced polymerization, including the evolution of molecular weight and dispersity \cite{photo2026}. Studies of cyclic polylactide and related cyclic polymers place diffusion measurements in an architecture-sensitive characterization setting \cite{cyclic2021,cyclic2023}. Low-molecular-weight chitosan provides a further materials application in which molecular weight matters \cite{chitosan2026}; that application is cited as motivation, not as evidence of ILT recovery or of performance by the present algorithms. On the inversion side, diffGA uses boxcar representations in a genetic-algorithm search for polymer mixtures \cite{diffga}, while dART adopts algebraic reconstruction for diffusion decays \cite{dart}. The regularized Kaczmarz work extends this algebraic perspective \cite{kaczmarz}. The TRAIn-MF reference also has a direct software lineage: Arrabal-Campos's \texttt{TRAIn\_DOSY\_MFV31} extends the original TRAIn to a multifrequency factorization \cite{train,mfv31}. The present reference is a later numerical revision of that extension, as detailed in Section~\ref{sec:algorithm}.

Finite exponential models introduce a second established lineage. Linear amplitudes can be eliminated from nonlinear least squares by variable projection \cite{golub}; constrained extensions require attention to active sets and differentiability \cite{shearer}. Nonnegative least squares (NNLS) supplies the conditional amplitude problem \cite{lawson}. Continuous sparse-measure methods, including sliding Frank--Wolfe, show that off-grid Laplace support recovery has its own mathematical theory under explicit hypotheses \cite{denoyelle}. The method considered here is a finite candidate search with support selection. It is not a new proof of unrestricted super-resolution and does not inherit guarantees from a different estimator.

Our focus is the interface between representation, optimization, and selection. We use a positive-measure formulation to compare continuous shared profiles with shared atomic rates. We derive the exact local residual derivative on a stable NNLS active set, explain the information in residualized kernels, and specify a selector that uses spectral coherence without deleting components solely because their heights are small. The selector is applied to two existing research implementations: resolution-adaptive inversion in its component mode (RAI) and diffusion orthogonal moment exchange (DOME). Their different proposal mechanisms feed a common final estimator, denoted RAI-S and DOME-S. The suffix identifies spectral-support selection; it does not designate a new acquisition or a distinct physical model.

The mathematical statements have deliberately limited scope. We prove finite-sampling nonuniqueness for unrestricted positive measures and contrast it with exact uniqueness in a bounded-order noiseless atomic class. We establish a grid-conversion bound, conditional uniqueness of amplitude subproblems, and conditional Gaussian identities motivating selection. These are transparent consequences and adaptations of established principles, rather than claims of priority. The empirical contribution is an auditable implementation and a frozen comparison showing a reduction in spectral leakage, together with failures in weak and proportional mixtures. A distributional multifrequency factorization reference, named \emph{TRAIn-MF}, is retained for finite-width profiles. The suffix MF denotes multifrequency processing, consistent with the historical source header, while matrix factorization describes its mathematical form. Its factor count is not equated with a chemical species count.

The organization follows this distinction. Section~\ref{sec:related} positions the work among related inverse-problem methods. Section~\ref{sec:theory} gives the measure model and mathematical properties. Section~\ref{sec:algorithm} specifies optimization and automatic selection. Sections~\ref{sec:experiments} and \ref{sec:results} describe the simulations and complete primary results. The discussion identifies what remains unresolved, including unknown or correlated noise, structural misspecification, and experimental validation.
